\documentclass[10pt]{article}

\usepackage[margin=1in]{geometry}
\usepackage{amsmath,amssymb}

\newcommand{\sg}{\operatorname{sg}}
\newcommand{\E}{\mathbb{E}}
\newcommand{\Dunary}{\mathcal{D}_{\mathrm{u}}}
\newcommand{\Danchor}{\mathcal{D}_{\mathrm{a}}}
\newcommand{\Lcfm}{\mathcal{L}_{\mathrm{CFM}}}
\newcommand{\Lkto}{\mathcal{L}_{\mathrm{KTO}}}

\title{Flow-KTO: Learning Flow-Based Robot Policies from Unary Feedback}
\author{Anonymous Authors}
\date{}

\begin{document}
\maketitle

\begin{abstract}
% Deferred by the confirmed writing sequence. Write the Abstract only after the
% Method, Experiments, and Conclusion have stabilized.
\end{abstract}

\section{Introduction}

% Move 1: state the complete problem.
During real-world robot deployment, a policy produces an action trajectory from
the current observation, and executing that trajectory changes the physical
scene. Collecting an alternative action under exactly the same conditions would
require reproducing object poses, contact states, sensor readings, and timing.
Such reconstruction is generally unavailable after execution. The deployed
trajectory can nevertheless be inspected afterward and labeled as desirable or
undesirable. The resulting deployment feedback contains one executed trajectory
and one unary judgment for each encountered context, but no same-context
counterfactual. It identifies behavior to encourage or suppress without
providing the alternative needed for a direct comparison. Reusing this feedback
is attractive precisely because it avoids recollecting counterfactual actions
or rewarded candidate groups on real hardware. The challenge is to turn this
direction-only deployment feedback into a valid policy update.

Flow-based robot policies add a separate optimization constraint. They generate
continuous action trajectories by integrating a learned vector field from noise
through multiple steps \cite{lipman2023flowmatching,chi2023diffusionpolicy}.
Their networks therefore do not emit terminal action likelihoods as direct
outputs. Naively differentiating a trajectory-level objective through the full
generation process can also yield unstable gradients when adapting such a
deployed policy. The problem is consequently defined by both the deployment
feedback and the policy interface: can a pretrained flow-based robot
policy be improved directly from unary feedback collected during real-world
deployment, without collecting additional robot interactions for optimization
or learning a reward or value critic?

% Move 2: review existing methods and identify the mismatch.
Existing post-training methods obtain a learning signal by adding comparison or
value structure. Direct Preference Optimization requires a preferred and a
dispreferred output for the same context \cite{rafailov2023dpo}; Proximal Policy
Optimization and Group Relative Policy Optimization instead form advantages
from rewarded current-policy trajectories or groups of candidates
\cite{schulman2017ppo,schulman2016gae,shao2024deepseekmath}. Recent flow-policy
methods illustrate critic-based alternatives. RECAP converts deployment
outcomes into binarized advantage conditioning through a learned value function,
using demonstrations, autonomous rollouts, and human corrections
\cite{physicalintelligence2025pi06star}. Q-learning with Adjoint Matching learns
a temporal-difference critic from reward-labeled transitions and uses adjoint
matching to turn its action gradients into supervision for the flow field
\cite{li2026qam}. These routes are viable under their respective data
interfaces, but they do not define the restricted update studied here, whose
only quality signal is a fixed set of individually judged deployment
trajectories. They first construct pairwise or group comparisons, collect
current-policy experience, or learn critic-derived targets. We instead ask
whether unary judgments can reach the flow policy without these intermediaries.

% Move 3: state the proposed response.
Without a learned intermediary, a unary judgment must be converted into a
comparison using quantities available from the policy itself. The first
comparison asks how the current policy's fit to a labeled deployment trajectory
has changed relative to the frozen pretrained policy. This reveals how training
has moved the recorded behavior, but not how that change stands relative to the
current policy's own outputs. A second, policy-centered comparison must therefore
come from trajectories generated by the current policy at the stored deployment
contexts. Recorded deployment actions cannot provide this reference because
they represent the behavior-data distribution; actions borrowed from other
contexts would additionally confound policy change with context--action
mismatch. Together, the two comparisons anchor each update to the pretrained
policy while interpreting it against the policy being optimized.

We propose Flow-KTO to implement this principle. Flow-KTO first measures the
change in conditional flow-matching fit for each labeled deployment trajectory
relative to a frozen reference policy. It then samples trajectories internally
from the current policy at stored deployment contexts, evaluates each sample
under the current and reference policies, and aggregates those comparisons into
a dynamic reference point. This construction keeps every labeled comparison
anchored to the pretrained policy while locating it within the current policy's
output distribution. The unary label determines which side of this
policy-centered reference should receive utility, and the exact bounded KTO
objective converts that direction into a policy update
\cite{ethayarajh2024kto}. The internally sampled trajectories are not assigned
quality labels and are never executed on the robot; they are used only to
estimate the current policy's reference point. Because
flow-matching differences need not share the numerical scale of explicit
likelihood ratios, Flow-KTO calibrates the inverse temperature to the scale of
these comparisons while keeping class-utility and objective-mixing weights
separate.

% Add the controlled-results summary only after the Experiments section fixes
% the matched baselines, estimator ablations, uncertainty, and robot protocol.
% Move 4: itemize the core contributions.
The paper makes three technical contributions:

\begin{itemize}
    \item We formulate direct post-training of a pretrained flow-based robot
    policy from a fixed collection of feedback gathered during real-world robot
    deployment, with one unary-labeled trajectory per encountered context and no
    requirement for paired alternatives, rewarded groups, additional robot
    interaction, or a learned reward or value critic.
    \item We introduce a policy-centered reference estimator in the native
    flow-matching training space. It uses internally sampled current-policy
    trajectories at stored deployment contexts and is distinct from references
    computed on recorded deployment actions or context-mismatched actions.
    \item We optimize the resulting label-directed comparison with exact
    bounded KTO and calibrate its inverse temperature to the scale of the
    flow-matching comparisons, separate from class-utility and objective-mixing
    weights.
\end{itemize}


\section{Method}

\subsection{Fixed Unary Robot Logs}

Let $x\in\mathcal{X}$ denote the context available to a robot policy, including
the observation history and any task conditioning, and let
$y\in\mathbb{R}^{H\times d_a}$ be an action chunk of horizon $H$. The unary
post-training set is
\begin{equation}
    \Dunary=\{(x_i,y_i,s_i)\}_{i=1}^{N},
    \qquad s_i\in\{+1,-1\},
    \label{eq:dataset}
\end{equation}
where $s_i=+1$ marks a desirable trajectory and $s_i=-1$ an undesirable one.
There is at most one logged trajectory for a context in the studied corpus; no
paired counterfactual $y'_i$ is assumed.

We start from a pretrained flow policy $\pi_{\mathrm{ref}}$ and initialize the
trainable policy $\pi_\theta$ from the same parameters. The reference remains
frozen. Post-training does not query an environment, and internally sampled
actions have no reward or quality label. We also allow an optional anchor set
$\Danchor$ of ordinary behavior-cloning trajectories. It is used only through
the original CFM training objective and is kept separate from the unary
feedback in Eq.~\eqref{eq:dataset}.

\subsection{A Reference-Relative CFM Improvement Score}

We first define how a complete action chunk is scored by the flow policy. Let
$\epsilon\sim p_0$ be a base-noise sample and let $t\sim p(t)$ denote a flow
time. A conditional probability path can be written as
\begin{equation}
    y_t=\alpha_t y+\sigma_t\epsilon,
    \qquad
    u_t(y,\epsilon)=\dot{\alpha}_t y+\dot{\sigma}_t\epsilon,
    \label{eq:path}
\end{equation}
where $u_t$ is the conditional target velocity. This notation includes the
linear optimal-transport path as a special case. For a vector field $v_j$ with
$j\in\{\theta,\mathrm{ref}\}$, we use the same schedule, action mask $M$, and
reduction as the pretrained policy:
\begin{align}
    \ell_j(t,\epsilon;x,y)
    &=\frac{w(t)}{\lVert M\rVert_1}
      \left\lVert
      M\odot\bigl(v_j(y_t,t\mid x)-u_t(y,\epsilon)\bigr)
      \right\rVert_2^2, \\
    \Lcfm^{j}(y\mid x)
    &=\E_{t,\epsilon}\left[\ell_j(t,\epsilon;x,y)\right].
    \label{eq:cfm}
\end{align}
All current/reference comparisons share the same draws $(t,\epsilon)$, masks,
preprocessing, and normalization. Pairing removes nuisance variation from the
difference without changing its expectation.

The Flow-KTO score is
\begin{equation}
    \Delta_\theta(x,y)
    =\Lcfm^{\mathrm{ref}}(y\mid x)
     -\Lcfm^{\theta}(y\mid x).
    \label{eq:delta}
\end{equation}
A positive value means that the current vector field fits the scored trajectory
better than the frozen reference under the shared teacher-forced protocol. With
$R$ Monte Carlo draws, we estimate Eq.~\eqref{eq:delta} by
\begin{equation}
    \widehat{\Delta}^{(R)}_\theta(x,y)
    =\frac{1}{R}\sum_{r=1}^{R}
      \left[\ell_{\mathrm{ref}}(t_r,\epsilon_r;x,y)
      -\ell_\theta(t_r,\epsilon_r;x,y)\right].
    \label{eq:delta-mc}
\end{equation}

The score has a precise but limited likelihood interpretation. For weightings
under which the denoising objective satisfies
$\mathcal{L}^{w}_{j}(y\mid x)=-\mathrm{ELBO}_{j}(y\mid x)+C(x,y)$,
as used in the FPO derivation \cite{mcallister2025fpo},
\begin{align}
    \Delta^w_\theta(x,y)
    &=\mathrm{ELBO}_{\theta}(y\mid x)
      -\mathrm{ELBO}_{\mathrm{ref}}(y\mid x) \\
    &=\log\frac{\pi_\theta(y\mid x)}
                      {\pi_{\mathrm{ref}}(y\mid x)}
      -\left(G_\theta(x,y)-G_{\mathrm{ref}}(x,y)\right),
    \label{eq:elbo-gap}
\end{align}
where $G_j$ is the model-dependent gap between log likelihood and its ELBO.
Equation~\eqref{eq:elbo-gap} is why a CFM difference can carry policy-improvement
information, and also why we do not identify it with an exact log-likelihood
ratio. Throughout, $\Delta_\theta$ is a reference-relative CFM, or
ELBO-related, improvement score.

\subsection{A Policy-Centered Dynamic Surrogate Reference}

Whether a labeled score is large or small depends on how the current policy has
moved away from its reference. We therefore center it on current-policy actions.
Let $q_{\Dunary}(x)$ be the empirical context distribution. The target reference
quantity is
\begin{equation}
    b_\theta
    =\E_{x\sim q_{\Dunary}}
       \E_{\widetilde{y}\sim\pi_\theta(\cdot\mid x)}
       \left[\Delta_\theta(x,\widetilde{y})\right].
    \label{eq:policy-center}
\end{equation}
This expectation follows the actions that $\pi_\theta$ currently proposes. It
does not assign utility to those actions.

For a batch of $B$ logged contexts and $K$ samples per context, the practical
estimator is
\begin{equation}
    \widehat{b}_{\mathcal{B}}
    =\sg\!\left[
      \frac{1}{BK}\sum_{i=1}^{B}\sum_{k=1}^{K}
      \widehat{\Delta}^{(R)}_\theta
      \left(x_i,\widetilde{y}_{ik}\right)
      \right],
    \quad
    \widetilde{y}_{ik}\sim\pi_\theta(\cdot\mid x_i).
    \label{eq:bhat}
\end{equation}
Sampling is performed without gradients. The sampled action is then
teacher-forced under both policies using paired CFM randomness, and the entire
aggregate is detached. In distributed training, workers aggregate the numerator
and count before applying the stop-gradient. Larger $K$, larger $R$, or
aggregation across batches can reduce variance; these choices affect estimator
quality and compute but not the estimand in Eq.~\eqref{eq:policy-center}.

Two superficially similar quantities answer different questions. A dataset
center
\begin{equation}
    b_{\mathrm{data}}
    =\E_{(x,y)\sim\Dunary}\left[\Delta_\theta(x,y)\right]
\end{equation}
measures relative fit on logged actions. A mismatched center such as
$\E_i[\Delta_\theta(x_i,y_{\rho(i)})]$, for a nontrivial permutation $\rho$,
measures response to context-action mismatch and may score out-of-distribution
actions. Both are useful diagnostics. Neither is interchangeable with the
current-policy expectation in Eq.~\eqref{eq:policy-center}. Because
$\Delta_\theta$ is a surrogate score, $b_\theta$ is likewise a dynamic
surrogate reference point, not an exact KL divergence.

\subsection{Exact KTO from a Signed Unary Margin}

For a labeled example, define the signed policy-centered margin
\begin{equation}
    m_i=s_i\left(
        \widehat{\Delta}^{(R)}_\theta(x_i,y_i)
        -\widehat{b}_{\mathcal{B}}
        \right).
    \label{eq:margin}
\end{equation}
Let $\lambda_D$ and $\lambda_U$ be the desirable and undesirable utility
weights, and write $\lambda_{s_i}=\lambda_D$ for $s_i=+1$ and
$\lambda_{s_i}=\lambda_U$ for $s_i=-1$. The per-example value and loss are
\begin{align}
    v_i(\theta)
    &=\lambda_{s_i}\,\sigma(\beta m_i), \\
    \ell_{\mathrm{KTO},i}(\theta)
    &=\lambda_{s_i}-v_i(\theta)
      =\lambda_{s_i}\left[1-\sigma(\beta m_i)\right], \\
    \Lkto(\theta)
    &=\E_{(x_i,y_i,s_i)\sim\Dunary}
      \left[\ell_{\mathrm{KTO},i}(\theta)\right].
    \label{eq:kto}
\end{align}
This is the exact bounded KTO utility \cite{ethayarajh2024kto}. There is no
logarithm around the sigmoid. Written separately, the two cases are
\begin{align}
    \ell_D
    &=\lambda_D\left[
      1-\sigma\!\left(\beta(\Delta_\theta-b_\theta)\right)
      \right], \\
    \ell_U
    &=\lambda_U\left[
      1-\sigma\!\left(\beta(b_\theta-\Delta_\theta)\right)
      \right].
    \label{eq:kto-cases}
\end{align}

The sign convention gives the intended CFM update directly. Treating the
reference point as stopped-gradient,
\begin{equation}
    \frac{\partial \ell_{\mathrm{KTO},i}}
         {\partial \widehat{\Delta}_i}
    =-s_i\lambda_{s_i}\beta
      \sigma(\beta m_i)\bigl[1-\sigma(\beta m_i)\bigr].
    \label{eq:delta-gradient}
\end{equation}
Minimization therefore increases $\Delta_i$ for a desirable example, which
reduces its current-policy CFM loss relative to the reference. For an
undesirable example it decreases $\Delta_i$, increasing the current-policy CFM
loss on that trajectory. The multiplier in Eq.~\eqref{eq:delta-gradient}
vanishes on both saturated sides; an example that is far from the boundary
receives little update from this term.

\subsection{Matching KTO to the CFM Margin Scale}

Three coefficients play different roles. The class weights
$\lambda_D,\lambda_U$ set the bounded utility amplitudes and class balance. An
outer coefficient $\eta_{\mathrm{KTO}}$ controls how the unary objective is
mixed with an optional retention objective. The inverse temperature $\beta$
sets the width of the active margin region. Conflating these controls would
turn score calibration into a change in the total KTO weight.

For exact KTO, the magnitude of the margin derivative is
\begin{equation}
    \left|\frac{\partial\ell_{\mathrm{KTO},i}}
                    {\partial m_i}\right|
    =\lambda_{s_i}\beta\,
      \sigma(\beta m_i)\bigl[1-\sigma(\beta m_i)\bigr].
    \label{eq:margin-gradient}
\end{equation}
Its peak is $\lambda_{s_i}\beta/4$ at $m_i=0$, and its active band has width
on the order of $1/\beta$. In our flow-policy training, raw centered CFM margins
are typically on the $10^{-3}$ to $10^{-2}$ scale. With
$\beta=\mathcal{O}(10^{-2})$, a $10^{-3}$ margin gives
$|\beta m|=\mathcal{O}(10^{-5})$; the utility remains effectively fixed at its
midpoint. We instead keep $\lambda_D$, $\lambda_U$, and
$\eta_{\mathrm{KTO}}$ unchanged and use $\beta=\mathcal{O}(10^{2})$. The
sigmoid then resolves changes at the observed CFM scale, and its gradient
automatically attenuates as either signed margin saturates.

Changing $\beta$ does raise the peak boundary gradient even though the utility
amplitude stays fixed. The optimizer learning rate, gradient clipping, score
normalization, and sampling variance must therefore be controlled separately.
The scale-matching claim is about the argument of the exact KTO utility; it does
not imply that a larger $\beta$ alone guarantees a better policy.

\subsection{Training Objective and Procedure}

The core objective is Eq.~\eqref{eq:kto}. When retention data are available, we
optionally add the original flow-training loss
\begin{align}
    \mathcal{L}_{\mathrm{anchor}}(\theta)
    &=\E_{(x,y)\sim\Danchor}\left[\Lcfm^\theta(y\mid x)\right], \\
    \mathcal{L}_{\mathrm{total}}(\theta)
    &=\eta_{\mathrm{KTO}}\Lkto(\theta)
      +\eta_{\mathrm{a}}\mathcal{L}_{\mathrm{anchor}}(\theta).
    \label{eq:total}
\end{align}
Setting $\eta_{\mathrm{a}}=0$ gives the pure unary objective. The anchor is a
retention mechanism, not an additional source of preference or reward; its
effect must be evaluated independently.

One optimization step proceeds as follows:
\begin{enumerate}
    \item Draw a unary batch from $\Dunary$ and, if used, an anchor batch from
    $\Danchor$.
    \item At each logged context, sample $K$ action chunks from
    $\pi_\theta$ without gradient tracking and without querying the robot.
    \item Teacher-force every sampled chunk through $\pi_\theta$ and
    $\pi_{\mathrm{ref}}$ with paired flow times and noise. Aggregate
    Eq.~\eqref{eq:bhat} across the batch and workers, then detach it.
    \item Score the logged unary trajectories with Eq.~\eqref{eq:delta-mc},
    form their signed margins, and compute the exact KTO loss in
    Eq.~\eqref{eq:kto}.
    \item Add the optional anchor term and update $\theta$. The reference
    policy, policy samples, and dynamic reference receive no gradient.
\end{enumerate}

For $B$ contexts, $K$ policy samples, $S$ flow-solver steps, and $R$ paired
CFM draws, reference estimation adds $BKS$ no-gradient vector-field evaluations
for sampling and $2BKR$ teacher-forced current/reference evaluations, in
addition to scoring the labeled examples. No reward model, value model, or
backpropagation through the sampler is required. This cost is model computation,
not environment interaction.

Flow-KTO assumes a fixed reference policy, binary trajectory-level judgments,
and a consistent CFM scoring protocol whose numerical scale can be measured.
It does not recover an exact terminal policy likelihood, and its dynamic
reference is not an exact KL divergence. Nor does the formulation rule out
critic-based offline learning or on-policy flow optimization when their richer
signals are available. It addresses the narrower case established at the
start: direct post-training from fixed unary robot logs.

\bibliographystyle{unsrt}
\bibliography{references}

\end{document}
